% Part: proof-theory
% Chapter: normalization
% Section: translations

\documentclass[../../../include/open-logic-section]{subfiles}

\begin{document}

\olfileid{pt}{nor}{tr}

\olsection{Nerjemahaké Antarane \Log{N2i}-\usetoken{P}{proof} Normal lan \Log{G2i}}

!!^{proof}s ing \Log{N2i} bisa diterjemahaké dadi !!{proof}s ing
\Log{G2i}, lan kosok baliné. Wiwitané, panerjemahan !!{proof}s
\Log{G2i} kang tanpa cut ngasilaké \Log{N2i}-!!{proof}s normal.

Kanggo nyatakaké asil iki, kita perlu istilah kang nggampangaké
sesambungan antarané antésèden sekuèn ing~\Log{N2i} lan
ing~\Log{G2i}: himpunan~$\Gamma'$ kang isiné !!{formula}s mawa
label \emph{cocog karo} multihimpunan~$\Gamma$ kang isiné
!!{formula}s tanpa label yèn, kanggo saben !!{formula}~$!A$,
cacahé !!{element}s ing~$\Gamma'$ kang awujud $x : !A$ ora
ngluwihi multiplisitas~$!A$ ing~$\Gamma$.

\begin{cor}
Yèn ana bukti \Log{G2i} tanpa cut tumrap
$\Gamma \Sequent \Delta$, ana bukti \Log{N2i} normal~$\delta$
tumrap $\Gamma' \Sequent !C$. Ing kéné $!C$ iku rumus ing
$\Delta$ yèn ora kosong, lan $!C \ident \lfalse$ yèn
$\Delta$ kosong. Himpunan~$\Gamma'$ kang
label cocog karo multihimpunan~$\Gamma$: kanggo saben
!!{formula}~$!A$, cacahé !!{element}s ing~$\Gamma'$ kang awujud
$x : !A$ ora ngluwihi multiplisitas~$!A$ (cacahé salinan~$!A$)
ing~$\Gamma$.
\end{cor}

\begin{proof}
Iki katon saka bukti \olref[nat][tgi]{prop:G2i-to-N2i}:
aturan !!{introduction} ditambahaké ing pungkasan !!{proof}s,
déning aturan !!{elimination} mung ing pérangan ndhuwur
!!{proof}s. Mula, ora ana aturan !!{introduction} kang
diselehaké ing sadhuwuré aturan !!a{elimination}.
\end{proof}

Nanging, panerjemahan aturan~\Cut{} ing
\olref[nat][tgi]{prop:G2i-to-N2i} bisa ngasilaké cut ing
\Log{N2i}-!!{proof}~$\delta$. Yèn !!{proof} \Log{G2i} pungkasané
nganggo~\Cut{}, kanthi sub-!!{proof}s langsung $\pi_1$ lan $\pi_2$
tumrap $\Gamma_1 \Sequent !A$ lan $!A, \Gamma_2 \Sequent \Delta_2$,
~$\delta_1$, yaiku panerjemahan~$\pi_1$, ditempelaké ing aksioma
$x:!A \Sequent !A$ kang ana ing~$\delta_2$, yaiku panerjemahan~$\pi_2$.
Yèn $\delta_1$ pungkasané nganggo aturan !!a{introduction} kang
!!{formula} utamané~$!A$, lan aksioma $x:!A \Sequent !A$ ing~$\delta_2$
dadi premis utama aturan !!a{elimination}, asilé ngandhut cut
ing~$\delta$.

!!{proof}s ing \Log{N2i} uga bisa diterjemahaké dadi !!{proof}s
ing~\Log{G2i}. Ing bukti \olref[nat][tni]{prop:N2-to-G2},
panerjemahan \Elim{\land}, \Elim{\lif}, \Elim{\lnot}, lan
\Elim{\lforall} ngasilaké inferènsi~\Cut{} ing !!{proof}
\Log{G2i}. Nanging, yèn \Log{N2i}-!!{proof} asalé normal,
inferènsi cut iki bisa diindhari.

Dhaptar kedadéan sekuèn $S_1$, \dots,~$S_n$ ing
\Log{N2i}-!!{proof}~$\delta$ diarani \emph{cabang} yèn $S_1$
minangka aksioma lan, saben $S_i$ dadi premis utama sawijiné
inferènsi, $S_{i+1}$ minangka konklusiné. Tegesé, cabang
nglacak sekuèn saka aksioma mudhun ing~$\delta$, nanging mandheg
ing premis tambahan aturan !!{elimination}. \emph{Cabang utama}
ing~$\delta$ yaiku cabang kang $S_n$ dadi sekuèn pungkasan.
Saben !!{proof}~$\delta$ nduwèni paling ora siji cabang utama,
amarga saben aturan nduwèni paling ora siji premis utama
(mung \Intro{\land} kang nduwèni loro).


\begin{prop}\ollabel{prop:normal-N2i-to-G2i} Yèn $\delta$
minangka $\Log{N2c}$- utawa $\Log{N2i}$-!!{proof} normal tumrap
$\Gamma \Sequent !A$, ana $\Log{G2c}$-!!{proof} (utawa
$\Log{G2i}$-!!{proof})~$\pi$ tanpa cut tumrap
$\Gamma' \Sequent \Delta$. Ing kéné $\Gamma'$ iku multihimpunan
!!{formula}s kang diasilaké saka~$\Gamma$ kanthi mbusak label;
$\Delta = \{!A\}$ yèn $!A$ dudu~$\lfalse$, lan
$\Delta = \emptyset$ yèn mangkono.
\end{prop}

\begin{proof}
Saiki induksi ditindakaké miturut cacahé inferènsi ing
!!{proof}~$\delta$ tumrap $\Gamma \Sequent !A$. Yèn ora ana
inferènsi ing~$\delta$, bukti iku aksioma
$x:!A \Sequent !A$; $\pi$ yaiku aksioma kang cocog,
$!A \Sequent !A$, utawa $\lfalse \Sequent \ $ yèn
$!A \ident \lfalse$.

Yèn $\delta$ pungkasané nganggo aturan inferènsi, kita mbedakaké kasus:
\begin{enumerate}
  \item Yèn $R$ aturan !!a{introduction} utawa~\FalseInt,
  panerjemahané padha karo ing bukti
  \olref[nat][tni]{prop:N2-to-G2}, lan ora mbutuhaké~\Cut.

  Yèn aturan pungkasan~$R$ ing~$\delta$ iku aturan
  !!a{elimination}, gatèkna prakara-prakara iki:
  \begin{enumerate}
    \item Cabang utama~$\delta$ ora ngliwati aturan
    !!{introduction}; yèn ngliwati, ana aturan
    !!a{introduction} utawa~\FalseInt{} kang disusul aturan
    !!a{elimination} ing cabang utama, mula $\delta$ ora normal.
    \item Amarga ora ana aturan !!{introduction} ing cabang-cabang
    utama~$\delta$, cabang utamané mung siji. (Mung
    \Intro{\land} kang nduwèni loro premis utama; mula cabang
    utama kang luwih saka siji mesthi ngliwati premis-premis
    aturan \Intro{\land}.)
    \item Yèn cabang utama ngemot premis utama \Elim{\lor}
    utawa \Elim{\lexists}, aturan kaya mangkono mung siji lan
    dadi aturan pungkasan, yaiku~$R$. (Yèn ora, cabang utama
    ngemot aturan \Elim{\lor} utawa \Elim{\lexists} kang
    konklusiné dadi premis utama aturan !!a{elimination};
    bukti iku ora normal amarga konversi permutasi bisa
    ditindakaké.)
  \item Ora ana aturan !!{elimination} saliyané \Elim{\lor}
  lan \Elim{\lexists} kang mbusak asumsi, yaiku ngowahi
  kontèks kiwa sekuèn. \Elim{\lor} lan \Elim{\lexists}
  mung mbusak asumsi ing
  sadhuwuré premis tambahan, mula mung ngowahi kontèks kiwa
  premis tambahané. Tegesé, yèn sekuèn $S_i$ ing cabang
  utama~$\delta$ nduwèni kontèks~$\Gamma_1$, konklusi~$S_{i+1}$
  saka inferènsi kang premis utamané~$S_i$ nduwèni kontèks
  $\Gamma_1' \supseteq \Gamma_1$.
  \item Mula, yèn $x: !C \Sequent !C$ minangka sekuèn
  paling dhuwur~$S_1$ ing cabang utama, $x:!C \in \Gamma$.
  \end{enumerate}
  Saiki kasus-kasus dibédakaké manut wujudé~$!C$. Amarga saben
  $S_i$ ing cabang utama dadi premis utama aturan
  !!a{elimination}, iki uga lumaku tumrap $x:!C \Sequent !C$.

  \item $!C \ident !D \land !E$. Mula $\delta$ awujud
  \[
  \Axiom$x:!D \land !E \fCenter !D \land !E$
  \RightLabel{\Elim{\land}[1]}
  \UnaryInf$x:!D \land !E \fCenter !D$
  \Deduce$x:!D \land !E, \Gamma_1 \fCenter !A$
  \DisplayProof
  \]
  Cabang utama kang ngemot $x:!D \land !E \Sequent !D$ ora
  ngliwati premis utama \Intro{\lif} utawa \Intro{\lnot},
  utawa premis tambahan \Elim{\lor} utawa \Elim{\lexists}.
  Mula !!{formula}s kontèks $x:!D \land !E$ ora owah
  sadawané cabang utama. Iki nerangaké sebabé kontèks
  sekuèn pungkasan kudu ngemot~$x:!D \land !E$. Salajengé,
  !!{proof}~$\delta$ bisa diowahi dadi !!{proof}~$\delta_1$:
  sub-!!{proof} kang pungkasané nganggo \Elim{\land}
  diganti $x:!D \Sequent !D$, lan $x:!D \land !E$ ing
  kontèks saben sekuèn sadawané cabang utama diganti
  $x:!D$. Dadi, kita ngowahi
  \[
  \bottomAlignProof
  \Axiom$x:!D \land !E \fCenter !D \land !E$
  \RightLabel{\Elim{\land}[1]}
  \UnaryInf$x:!D \land !E \fCenter !D$
  \Deduce$x:!D \land !E, \Gamma_1 \fCenter !A$
  \DisplayProof
  \quad\text{ dadi }\quad
  \bottomAlignProof
  \Axiom$x:!D \fCenter !D$
  \Deduce$x:!D, \Gamma_1 \fCenter !A$
  \DisplayProof
  \]
  Hipotèsis induksi lumaku tumrap~$\delta_1$, amarga
  cacahé inferènsi ing~$\delta$ padha karo cacahé
  inferènsi ing~$\delta_1$ ditambah~$1$.

  Miturut hipotèsis induksi, ana \Log{G2i}-!!{proof}~$\pi_1$
  tumrap $!D, \Gamma_1' \Sequent \Delta$. Kita olèh
  !!{proof}~$\pi$ kanthi ngetrapaké \LeftR{\land}:
  \[
  \AxiomC{}
  \RightLabel{$\pi_1$}
  \Deduce$!D, \Gamma_1' \fCenter !A$
  \RightLabel{\LeftR{\land}}
  \UnaryInf$!D \land !E, \Gamma_1' \fCenter !A$
  \DisplayProof
  \]
  Yèn $!A \ident \lfalse$, suksedèné tetep kosong,
  tanpa \RightR{\Weakening}; tuladhané,
  \[
  \AxiomC{}
  \RightLabel{$\pi_1$}
  \Deduce$!D, \Gamma_1' \fCenter $
  \RightLabel{\LeftR{\land}}
  \UnaryInf$!D \land !E, \Gamma_1' \fCenter $
  \DisplayProof
  \]
  \item Yèn $!C \ident !D \lor !E$, iku
  mesthi dadi premis utama \Elim{\lor}, lan inferènsi iki
  mesthi inferènsi pungkasan~$R$ ing cabang utama.
  Tegesé, $\delta$ awujud:
  \[
  \Axiom$x:!D \lor !E \fCenter !D \lor !E$
  \AxiomC{}
  \RightLabel{$\delta_1$}
  \Deduce$y:!D, \Gamma_1 \fCenter!A$
  \AxiomC{}
  \RightLabel{$\delta_2$}
  \Deduce$z:!E, \Gamma_2 \fCenter !A$
  \RightLabel{\Elim{\lor}}
  \TrinaryInf$x:!D \lor !E, \Gamma_1, \Gamma_2 \fCenter !A$
  \DisplayProof
  \]
  Hipotèsis induksi lumaku tumrap !!{proof}s $\delta_1$ lan
  $\delta_2$; kita olèh \Log{G2i}-!!{proof}s $\pi_1$ lan $\pi_2$
  tumrap $!D, \Gamma_1' \Sequent \Delta$ lan
  $!E, \Gamma_2' \Sequent \Delta$. Kanthi ngetrapaké
  $\LeftR{\lor}$, kita olèh~$\pi$:
  \[
  \AxiomC{}
  \RightLabel{$\pi_1$}
  \Deduce$!D, \Gamma_1' \fCenter !A$
  \AxiomC{}
  \RightLabel{$\pi_2$}
  \Deduce$!E, \Gamma_2' \fCenter !A$
  \RightLabel{\LeftR{\lor}}
  \BinaryInf$!D \lor !E, \Gamma_1', \Gamma_2' \fCenter !A$
  \DisplayProof
  \]
  Yèn $\Elim{\lor}$ ora mbusak $y: !D$ utawa $z: !E$
  (yaiku yèn salah sijiné ora ana ing kontèks premis
  tambahan), \LeftR{\Weakening} ditambahaké miturut
  kabutuhan. Yèn $!A \ident \lfalse$, suksedèné tetep
  kosong, tanpa \RightR{\Weakening}. Tuladhané,
  \[
  \AxiomC{}
  \RightLabel{$\pi_1$}
  \Deduce$\Gamma_1' \fCenter $
  \RightLabel{\LeftR{\Weakening}}
  \UnaryInf$!D, \Gamma_1' \fCenter $
  \AxiomC{}
  \RightLabel{$\pi_2$}
  \Deduce$\Gamma_2' \fCenter $
  \RightLabel{\LeftR{\Weakening}}
  \UnaryInf$!E, \Gamma_2' \fCenter $
  \RightLabel{\LeftR{\lor}}
  \BinaryInf$!D \lor !E, \Gamma_1', \Gamma_2' \fCenter $
  \DisplayProof
  \]
  
  \item $!C \ident !D \lif !E$. Mula $\delta$ awujud
  \[
  \Axiom$x:!D \lif !E \fCenter !D \lif !E$
  \AxiomC{}
  \RightLabel{$\delta_1$}
  \Deduce$\Gamma_1 \fCenter!D$
  \RightLabel{\Elim{\lif}}
  \BinaryInf$x:!D \lif !E, \Gamma_1 \fCenter !E$
  \RightLabel{$\delta_2$}
  \Deduce$x:!D \lif !E, \Gamma_1, \Gamma_2 \fCenter !A$
  \DisplayProof
  \]
  Cabang utama kang ngemot $x:!D \lif !E,
  \Gamma_1 \Sequent !E$ ora ngliwati premis utama
  \Intro{\lif} utawa \Intro{\lnot}, utawa premis tambahan
  \Elim{\lor} utawa \Elim{\lexists}. Mula !!{formula}s
  kontèks ing $x:!D \lif !E, \Gamma_1$ ora owah
  sadawané cabang utama. Iki nerangaké sebabé kontèks
  sekuèn pungkasan kudu ngemot~$x:!D \lif !E, \Gamma_1$.
  Salajengé, fragmèn !!{proof}~$\delta_2$ bisa diowahi
  dadi !!{proof}~$\delta_2'$ kang bener: sub-!!{proof}
  kang pungkasané nganggo \Elim{\lif} diganti
  $x:!E \Sequent !E$, lan $x:!D \lif !E, \Gamma_1$
  ing kontèks saben sekuèn sadawané cabang utama diganti
  $x:!E$. Dadi, kita ngowahi
  \[
  \bottomAlignProof
  \Axiom$x:!D \lif !E, \Gamma_1 \fCenter !E$
  \RightLabel{$\delta_2$}
  \Deduce$x:!D \lif !E, \Gamma_1, \Gamma_2 \fCenter !A$
  \DisplayProof
  \quad\text{ dadi }\quad
  \bottomAlignProof
  \Axiom$x:!E \fCenter !E$
  \RightLabel{$\delta_2'$}
  \Deduce$x:!E, \Gamma_2 \fCenter !A$
  \DisplayProof
  \]
  Hipotèsis induksi lumaku tumrap $\delta_1$ lan $\delta_2'$,
  amarga cacahé inferènsi ing~$\delta$ iku gunggungé
  cacah inferènsi ing~$\delta_1$ lan~$\delta_2$, ditambah~$1$
  kanggo inferènsi \Elim{\lif} ing wiwitan cabang utama.
  (Yèn inferènsi iku uga inferènsi pungkasan~$R$ ing~$\delta$,
  $\delta_1$ ngemot $0$ inferènsi.)

  Miturut hipotèsis induksi, ana \Log{G2i}-!!{proof}s $\pi_1$
  tumrap $\Gamma_1' \Sequent \Delta_D$ lan~$\pi_2$
  tumrap $!E, \Gamma_2' \Sequent \Delta_A$. Ing kéné
  $\Delta_D$ kosong yèn $!D \ident \lfalse$ lan awujud
  $\{!D\}$ yèn ora; $\Delta_A$ ditemtokaké kanthi cara
  kang padha saka~$!A$. Kita olèh !!{proof}~$\pi$
  kanthi ngetrapaké \LeftR{\lif}:
  \[
  \AxiomC{}
  \RightLabel{$\pi_1$}
  \Deduce$\Gamma_1' \fCenter !D$
  \AxiomC{}
  \RightLabel{$\pi_2$}
  \Deduce$!E, \Gamma_2' \fCenter !A$
  \RightLabel{\LeftR{\lif}}
  \BinaryInf$!D \lif !E, \Gamma_1', \Gamma_2' \fCenter !A$
  \DisplayProof
  \]
  Yèn $!D \ident \lfalse$, kita nambah
  \RightR{\Weakening} mung ing premis kapisan supaya
  aturan implikasi kiwa bisa ditrapaké. Yèn $!A \ident \lfalse$,
  suksedèn ing premis kapindho lan konklusi tetep kosong;
  ora perlu pelemahan tengen ing konklusi. Tuladhané,
  \[
  \AxiomC{}
  \RightLabel{$\pi_1$}
  \Deduce$\Gamma_1' \fCenter $
  \RightLabel{\RightR{\Weakening}}
  \UnaryInf$\Gamma_1' \fCenter !D$
  \AxiomC{}
  \RightLabel{$\pi_2$}
  \Deduce$!E, \Gamma_2' \fCenter $
  \RightLabel{\LeftR{\lif}}
  \BinaryInf$!D \lif !E, \Gamma_1', \Gamma_2' \fCenter $
  \DisplayProof
  \]

\end{enumerate}
Kasus-kasus liyané ditangani kanthi cara kang padha lan
dipasrahaké dadi latihan.
\end{proof}

\begin{cor}
Saben !!{formula} kang katon ing bukti \Log{N1i} utawa
\Log{N2i} normal iku sub-!!{formula} saka
!!{formula} pungkasan utawa saka asumsi-asumsi kang isih
mbukak ing sekuèn pungkasan.
\end{cor}

\begin{proof}
  Miturut \olref{prop:normal-N2i-to-G2i}, saben
  \Log{N2i}-!!{proof} normal bisa diterjemahaké dadi
  \Log{G2i}-!!{proof}. Saka pamriksan bukti katon yèn
  saben !!{formula} ing \Log{N2i}-!!{proof} uga katon
  ing \Log{G2i}-!!{proof}. \Log{G2i}, kang ora nganggo
  aturan \Cut{}, nduwèni sipat sub-!!{formula}.
\end{proof}

\end{document}
